\section{Harnes, Monitors, and Policies}\label{sec:method}

\subsection{Episode loop and label semantics}

An episode is a (task, fault, monitor, policy) tuple. The harness steps the reference agent; at each
step the environment returns an observation, the monitor scores it, and any alarm is handed to the
policy, which may enqueue a \emph{recovery program}---a short sequence of operator steps executed
before the agent's own next step. Recovery therefore costs real steps and real tool calls; it is
never a free relabelling of the transcript.

Every step carries: the tool call, its arguments, expected state, observed state, whether they
matched, a vector of monitor features, token counts, latency, and a USD cost under a versioned price
sheet (\texttt{arb-price-2026.11}) with deterministic jitter. A trajectory additionally carries the
fault plan (class, injection step), the detection point, the recovery action taken, the number of
retries, and the verified outcome.

Two label conventions matter and are easy to get wrong:

\begin{itemize}
\item \textbf{Detection is only counted from the injection step onward.} An alarm that fires before
the fault exists is a false alarm, not a lucky early detection. In practice this turned the
``time-to-detection'' distribution from contaminated to clean (Fig.~\ref{fig:ttd}), and the residual
early-alarm rate on the harmful set is \EarlyAlarmRate\%.
\item \textbf{Success is read from world state.} The agent's own claim is recorded as a separate
field and is never consulted by the grader. This is what makes the self-report gap in
Sec.~\ref{sec:results} measurable.
\end{itemize}

\subsection{Monitor features}

Nineteen features are computed per step, grouped into families: error codes and argument-schema
violations; version staleness (observed version lag, stale flag); cross-source disagreement
(conflict flag, disagreement count); communication (message lost, sender mismatch, unacknowledged);
policy conformance (mutation without a prior read of the mutated field, re-read violations);
progress (no-progress runs, repeated identical calls, retry streak); and cost pressure
(tokens and latency relative to the task's expectation). Features are computed from the observable
trace only: no feature reads the fault plan, the injection step, the grader verdict, or expected
state.

The supervisor is evaluated under two regimes. \emph{Instrumented} exposes the harness telemetry
flags (staleness, conflict, message-loss). \emph{Blackbox} zeroes them, leaving only error strings,
payloads, version counters and timing---i.e.\ what an external observer with no cooperation from
the runtime would see. The gap between regimes is the \emph{detectability gap}, and it is the honest
measure of how much a deployment can expect from a monitor that cannot instrument its own tools.

\subsection{Detectors}

\paragraph{Rule detector.} Thirteen rules over the feature families, each mapping a conjunction of
signals to a class with a confidence; the step score is the maximum over rules. Rules are
hand-written from the taxonomy, not tuned on the test splits.

\paragraph{Supervised detector.} A multinomial logistic regression over the 19 features (8 classes
including \emph{none}), trained in numpy with class balancing and L2. It is deliberately small: with
a few thousand labelled steps we are interested in whether the feature set carries signal, not in
winning a modelling contest. The decision threshold is swept on \emph{dev} to maximise F1
(selected $\tau=\SupThreshold$), and calibration is reported as ECE and Brier score since a monitor
that is badly calibrated will flood the policy with false alarms.

\subsection{Recovery policies}

\begin{itemize}
\item \textbf{No recovery.} Detector present, alarms ignored. This is the ``aggressive agent''
baseline and defines harm.
\item \textbf{Retry only.} On tool error, re-issue the same call up to a budget. This is the
degenerate self-correction loop, and it is the baseline that Reflexion-style methods must beat.
\item \textbf{Rule-based.} Detected class $\rightarrow$ canonical operator program
(Table~\ref{tab:taxonomy}), with guards: no more than one recovery per detected event, a repeated
-failure guard that escalates after a second identical failure, and no retries after the budget.
\item \textbf{Model-based.} A learned \emph{recovery action ranker} chooses between the canonical
repair for the detected class and a clean containment. The ranker is a set of per-operator logistic
value models, $P(\text{success}\mid a, x)=\sigma(w_a^\top x + b_a)$, over the posterior class
distribution, the alarm score, step pressure, error/staleness/conflict flags, retry streak and
retries used. Per-operator heads are necessary: a single model over action one-hots has no
action$\times$feature interaction, so its argmax is constant across all inputs and it degenerates to
``which action is best on average''. Operators with fewer than 25 logged executions are not
rankable and are excluded from the candidate set rather than being scored from a prior.
\end{itemize}

The ranker is fit on \emph{deliberate} decisions only---logged recovery choices from the rule-based
policy on the train split. Blind retries are excluded because they are not a choice of operator, and
including them teaches the model ``retry is bad'' as a global prior. This is off-policy learning
from a single logging policy, and it is the main reason we expect the learned policy to underperform
(Sec.~\ref{sec:results}); we report it as a boundary result rather than tuning it into a win.

\subsection{Metrics, intervals, and cost}

\paragraph{Detection.} Step-level labels (is a fault active at this step) give precision, recall, F1
and threshold-free AUC; class accuracy among positive steps measures whether the monitor knows
\emph{which} failure it is. Event-level recall asks whether the fault was ever detected at all.

\paragraph{Recovery.} Recovery rate conditional on class, plus attempt rate (did the policy try),
containment rate (for unrecoverable faults: no successful mutation after the fault), and
time-to-detection (\ttd) and time-to-recovery (\ttr) in steps.

\paragraph{Propagation.} Corrupted-step count (steps whose observed state diverges from expected)
and downstream-impact count (steps that were still wrong after the fault window), reported per
policy and per class.

\paragraph{Cost.}\label{sec:cost} Tool calls, tokens, wall-clock latency and USD per episode, plus USD per
\emph{recovered} episode---the amortised price of a repair, which is the number a deployment
actually budgets.

\paragraph{Intervals.} All intervals are 95\% cluster bootstrap \citep{efron1993bootstrap} resampling
\emph{task templates}, not episodes: episodes sharing a template share the procedure and the fault
schedule shape, so treating them as independent would understate uncertainty. Policy comparisons use
a paired bootstrap over templates on the same episode set. Single-episode proportions use Wilson
intervals \citep{wilson1927}.
